\begin{tabularx}{\textwidth}{p{0.24\textwidth} p{0.35\textwidth} X}
\toprule
设定/参数 & 事实来源 & 模型中的限定 \\
\midrule
住院48小时全项目完成 & 赛题直接给定 & 事件硬截止时刻 \\
5\%特殊检查取上界时长 & 赛题直接给定比例；附件无个体标签 & 评价/预热分别按最近整数抽取5\% \\
项目时长区间与空腹10点窗口 & 赛题直接给定 & 名义/上界时长与硬时窗 \\
平峰$Q_{0.45}$--$Q_{0.55}$、高峰$Q_{0.90}$ & 运营分位阈值，非赛题指定日期 & 用真实日需求识别代表性容量水平 \\
22台设备能力、床旁7号机 & 附件表1与设备表文本证据 & 项目--设备兼容集 \\
50个项目采用类别级能力回退 & 项目名称与表1逐字/语义核心匹配失败时，仅沿用表1同项目族具体证据 & 主情景逐项留痕；严格证据情景禁用回退 \\
08:00--12:00、13:00--17:00日班模板 & 附件表6训练期首末报告活动重建 & 可复现容量模板，非未来正式排班 \\
星期--5分钟医生并发容量 & 训练日重构活动区间的同刻人数中位数 & 仅为总并发代理；不声称掌握未来实名班表或专项资质 \\
相邻报告间隔$(0,60]$分钟 & 扫查起止缺失下的局部代理阈值 & 只用于问题一相对吞吐证据，不进入主排程时长 \\
5分钟时间网格 & 计算离散化，附件无此字段 & 时长向上取整，不缩短检查 \\
跨室转运10分钟 & 附件无转运时间 & 主情景；5/15分钟边界检验 \\
憋尿准备60分钟 & 附件无时长；医院公开流程建议提前1小时\cite{hrhospital2025} & 主情景；45/90分钟边界检验 \\
门诊/体检报告日作为计划日 & 附件无历史预约日 & 回溯压力情景，不声称还原历史预约 \\
最后任务结束代理报告提交 & 附件无扫查起止与报告书写时长 & 官方KPI的乐观代理，明示适用边界 \\
2日预热与2日截止尾窗 & 由最大48小时截止期推出 & 只携入跨边界在制任务并覆盖3月31日开单截止 \\
\bottomrule
\end{tabularx}
